\begin{figure*}[b!]
    \begin{subfigure}{0.5\textwidth}
        \centering
        \includegraphics{img/poly_100M_pause_5_app.pdf}
        \label{fig:poly_5}
        \caption{5 pause tokens}
    \end{subfigure}
    \hfill
    \begin{subfigure}{0.5\textwidth}
        \centering
        \includegraphics{img/poly_100M_pause_10_app.pdf}
        \label{fig:poly_10}
        \caption{10 pause tokens}
    \end{subfigure}

    \caption{\textbf{Accuracy on a polynomial arithmetic task with varying number of operations per expression and pause tokens.} We train and evaluate models on the polynomial arithmetic task described in Section~\ref{sect:algorithmic}, modified by the addition of \emph{pause tokens} \citep{goyal2023think}: between the question and the equality sign that indicates the beginning of the answer, we add a constant number of pause tokens both in training and evaluation. For both a variant with five and with ten pause tokens, respectively, we observe comparable improvements from using multi-token prediction to the ones obtained in the case without pause tokens (Figure~\ref{fig:poly_0}).
    }
    \label{fig:poly_pause}
\end{figure*}
